% GENERATED FILE. Edit canonical lessons, then run npm run generate:books.
% canonical-source-hash: fnv1a64:ca255400a529108f
% canonical-lessons: FA-C115-shagerd, FA-C115-doktor, FA-C115-parastar, FA-C115-ashpaz, FA-C115-kargar

\chapter{Pupil, Doctor, Nurse, Cook, Worker}
\label{ch:fa-pupil-doctor-nurse-cook-worker}
\hlchaptermodality{115}

\begin{chapteropening}
\textbf{By the end of this chapter:} \emph{I can say a pupil, a doctor, a nurse, a cook and a worker.}

Complete the last lesson of chapter 115: I can say a pupil, a doctor, a nurse, a cook and a worker.
\end{chapteropening}

\section[\texorpdfstring{\fa{شاگرد} (shâgerd)}{شاگرد (shâgerd)}]{\fa{شاگرد} (shâgerd) — a pupil}

\label{lesson:FA-C115-shagerd}



\begin{quote}
Before the new one: say the Persian for an exam, then the Persian for a teacher.
\end{quote}

\subsection*{You'll want to know: \fa{شاگرد}}
\textbf{\fa{شاگرد}} — \emph{shâgerd} — \textquotedblleft{}a pupil\textquotedblright{}.

\begin{cousinweb}
\emph{shâgerd} is a native Persian word for a pupil or apprentice, but no source used here gives its older form, and its origin is not settled.
\end{cousinweb}

\begin{grammarlens}[title={What you've built}]
One more word for this chapter.
\end{grammarlens}

\subsection*{Your turn}
\begin{itemize}
\raggedright
  \item \emph{Say it:} \emph{shâgerd}
  \item \emph{Say it:} \emph{shâgerd}, once more
  \item \emph{Say it:} say \emph{shâgerd}
  \item \emph{Recall:} say the Persian for an exam, then the Persian for a teacher, then say \emph{shâgerd} again
\end{itemize}

\begin{tcolorbox}[breakable,colback=teal!4,colframe=teal!35!black,title={Before you move on}]
What does \fa{شاگرد} mean? (\textquotedblleft{}A pupil\textquotedblright{}.) Say it once more.
\end{tcolorbox}
\section[\texorpdfstring{\fa{دکتر} (doktor)}{دکتر (doktor)}]{\fa{دکتر} (doktor) — a doctor}

\label{lesson:FA-C115-doktor}



\begin{quote}
Before the new one: say the Persian for a teacher, then the Persian for a pupil.
\end{quote}

\subsection*{You'll want to know: \fa{دکتر}}
\textbf{\fa{دکتر}} — \emph{doktor} — \textquotedblleft{}a doctor\textquotedblright{}.

\begin{cousinweb}
\emph{doktor} is borrowed from French \emph{docteur}, \textquotedblleft{}doctor\textquotedblright{} (EtymDB).
\end{cousinweb}

\begin{grammarlens}[title={What you've built}]
One more word for this chapter.
\end{grammarlens}

\subsection*{Your turn}
\begin{itemize}
\raggedright
  \item \emph{Say it:} \emph{doktor}
  \item \emph{Say it:} \emph{doktor}, once more
  \item \emph{Say it:} say \emph{doktor}
  \item \emph{Recall:} say the Persian for a teacher, then the Persian for a pupil, then say \emph{doktor} again
\end{itemize}

\begin{tcolorbox}[breakable,colback=teal!4,colframe=teal!35!black,title={Before you move on}]
What does \fa{دکتر} mean? (\textquotedblleft{}A doctor\textquotedblright{}.) Say it once more.
\end{tcolorbox}
\section[\texorpdfstring{\fa{پرستار} (parastâr)}{پرستار (parastâr)}]{\fa{پرستار} (parastâr) — a nurse}

\label{lesson:FA-C115-parastar}



\begin{quote}
Before the new one: say the Persian for a pupil, then the Persian for a doctor.
\end{quote}

\subsection*{You'll want to know: \fa{پرستار}}
\textbf{\fa{پرستار}} — \emph{parastâr} — \textquotedblleft{}a nurse\textquotedblright{}.

\begin{cousinweb}
\emph{parastâr} is a native Persian word, inherited rather than borrowed: it continues Middle Persian \emph{parastār}, \textquotedblleft{}servant, attendant\textquotedblright{} (EtymDB). Persian now uses it for a nurse.
\end{cousinweb}

\begin{grammarlens}[title={What you've built}]
One more word for this chapter.
\end{grammarlens}

\subsection*{Your turn}
\begin{itemize}
\raggedright
  \item \emph{Say it:} \emph{parastâr}
  \item \emph{Say it:} \emph{parastâr}, once more
  \item \emph{Say it:} say \emph{parastâr}
  \item \emph{Recall:} say the Persian for a pupil, then the Persian for a doctor, then say \emph{parastâr} again
\end{itemize}

\begin{tcolorbox}[breakable,colback=teal!4,colframe=teal!35!black,title={Before you move on}]
What does \fa{پرستار} mean? (\textquotedblleft{}A nurse\textquotedblright{}.) Say it once more.
\end{tcolorbox}
\section[\texorpdfstring{\fa{آشپز} (âshpaz)}{آشپز (âshpaz)}]{\fa{آشپز} (âshpaz) — a cook}

\label{lesson:FA-C115-ashpaz}



\begin{quote}
Before the new one: say the Persian for a doctor, then the Persian for a nurse.
\end{quote}

\subsection*{You'll want to know: \fa{آشپز}}
\textbf{\fa{آشپز}} — \emph{âshpaz} — \textquotedblleft{}a cook\textquotedblright{}.

\begin{cousinweb}
\emph{âshpaz} is built from \emph{âsh}, \textquotedblleft{}soup, cooked food\textquotedblright{}, and \emph{-paz}, \textquotedblleft{}cooking\textquotedblright{}, the present stem of \emph{pokhtan}, \textquotedblleft{}to cook\textquotedblright{}: the one who cooks. \emph{pokhtan} is inherited, usually traced to Middle Persian \emph{puxtan}.
\end{cousinweb}

\begin{grammarlens}[title={What you've built}]
One more word for this chapter.
\end{grammarlens}

\subsection*{Your turn}
\begin{itemize}
\raggedright
  \item \emph{Say it:} \emph{âshpaz}
  \item \emph{Say it:} \emph{âshpaz}, once more
  \item \emph{Say it:} say \emph{âshpaz}
  \item \emph{Recall:} say the Persian for a doctor, then the Persian for a nurse, then say \emph{âshpaz} again
\end{itemize}

\begin{tcolorbox}[breakable,colback=teal!4,colframe=teal!35!black,title={Before you move on}]
What does \fa{آشپز} mean? (\textquotedblleft{}A cook\textquotedblright{}.) Say it once more.
\end{tcolorbox}
\section[\texorpdfstring{\fa{کارگر} (kârgar)}{کارگر (kârgar)}]{\fa{کارگر} (kârgar) — a worker}

\label{lesson:FA-C115-kargar}



\begin{quote}
Before the new one: say the Persian for a nurse, then the Persian for a cook.

Say \textbf{\fa{دوست} \fa{داشتن}}, then \textbf{\fa{کمک} \fa{کردن}}: two verbs, each in two pieces.
\end{quote}

\subsection*{You'll want to know: \fa{کارگر}}
\textbf{\fa{کارگر}} — \emph{kârgar} — \textquotedblleft{}a worker\textquotedblright{}.

\begin{cousinweb}
\emph{kârgar} is a native Persian word, inherited rather than borrowed: it continues Middle Persian \emph{kārgar}, \textquotedblleft{}worker\textquotedblright{} (EtymDB). It is \emph{kâr}, \textquotedblleft{}work\textquotedblright{}, with \emph{-gar}, \textquotedblleft{}doer\textquotedblright{}.
\end{cousinweb}

\begin{grammarlens}[title={What you've built}]
One more word for this chapter.
\end{grammarlens}

\subsection*{Your turn}
\begin{itemize}
\raggedright
  \item \emph{Say it:} \emph{kârgar}
  \item \emph{Say it:} \emph{kârgar}, once more
  \item \emph{Say it:} say \emph{kârgar}
  \item \emph{Recall:} say the Persian for a nurse, then the Persian for a cook, then say \emph{kârgar} again
\end{itemize}

\begin{tcolorbox}[breakable,colback=teal!4,colframe=teal!35!black,title={Before you move on}]
What does \fa{کارگر} mean? (\textquotedblleft{}A worker\textquotedblright{}.) Say it once more.
\end{tcolorbox}
